% Source: https://artofproblemsolving.com/community/c3h1611115_how_to_solve_19_in_the_patterns_stretch_of_the_mathcounts_handbook
% Stable identity: aops-topic-1611115-q19
Note that $f(x)+3=2x+6=2(x+3)$: applying $f$ doubles the distance from $-3$. Hence $f^n(x)+3=2^n(x+3)$, i.e.
\[
f^n(x)=2^nx+3\left(2^n-1\right).
\]
(Check: $n=2$ gives $4x+9$.) For $n=5$: $a=32$ and $b=3\cdot31=93$, so $a+b=\boxed{125}$.

(Directly: $f^3(x)=8x+21$, $f^4(x)=16x+45$, $f^5(x)=32x+93$. Also $a+b=f^5(1)$, and $1\to5\to13\to29\to61\to125$.)
